\section{Approach}
\label{sec:method}

\subsection{Overview}
Figure~\ref{fig:method-overview} presents the workflow of PyTT, our
implementation of interprocedural type-state analysis. Its inputs are a
Python repository, declared entries, dependency contracts, a Python
semantic configuration, and exploration budgets. The top row constructs
a call graph using pyflow~\cite{pyflow2026}, then analyzes functions in
callee-first order. Recursive functions are grouped into strongly
connected components; the topological order applies to the resulting
component graph. This call graph remains fixed throughout the analysis.

\begin{figure}[!htbp]
\centering
\includegraphics[width=\linewidth]{methodOverview.pdf}
\caption{Overview of PyTT. A fixed call graph orders function analysis.
For each function, Context-free Code Simplification supplies shared facts
to PSTC Generation. A PSTC contains entry assertions ($P$), ordered path
behavior, and exit assertions ($Q$). Type conflicts enter Type Conflict
Analysis, which expands relevant callee paths and re-infers the resulting
trace. Revised path contracts update callee PSTCs and dependent callers.}
\Description{The top row connects a repository to call graph generation,
a call graph, ordered function analysis, and bug reports. The bottom panel
shows function code and callee PSTCs entering three phases. The first
extracts and checks entry-independent code to supply type facts. The second
extracts type paths and resolves their type constraints to form a PSTC.
The third expands callee PSTCs and analyzes a conflict trace. Its outputs
are bug reports or a missing path contract that feeds back to update the
callee PSTC. A PSTC inset shows two entry, path, and exit cases, including
an attribute changing from type B to C.}
\label{fig:method-overview}
\end{figure}

The lower panel expands the analysis of one function into three phases.
Function code and available \emph{callee PSTCs} supply the local and
interprocedural information used in these phases. A \emph{path-sensitive
type contract} (PSTC) is our representation of a function's type-state
relations: each path contract connects entry assertions, ordered path
behavior, and exit assertions. We define it in Section~\ref{sec:state-model}.

\paragraph{Phase I: Context-free Code Simplification.}
The analyzer separates computations whose type behavior can be determined
without the unknown entry types and state from the remaining computations.
It checks the extracted part and supplies type facts at their uses in the
remaining code. In the figure, the available contract of \texttt{g} allows
\texttt{v = g(c)} with \texttt{c = 1} to supply the fact \texttt{v: int}
to the subsequent call \texttt{h(v, a)}. The extracted checks and effects
remain part of the function's behavior.

\paragraph{Phase II: PSTC Generation.}
\emph{Type Path Extraction} obtains explicit paths through the dependent
code. \emph{Type Constraint Resolution} combines their operations with
shared type facts and callee PSTCs, using the LLM to propose implicit type
alternatives and corresponding assertions. The analyzer assembles these
relations with the ordered checks and effects. Inferred relations are used
for propagation, and conflicts encountered during construction or
composition are passed to the third phase.

\paragraph{Phase III: Type Conflict Analysis.}
\emph{Callee PSTC Expanding} selects the callee type paths relevant to a
conflict and connects them to the caller's relevant behavior.
\emph{Conflict Path Analysis} re-infers this expanded trace in its calling
context. A confirmed incompatible behavior produces the corresponding
error report. A missing or inaccurate relation instead yields a new or
revised path contract, shown as \emph{Missing Path Contract} in the figure.
The \emph{Update Callee PSTC} feedback revises the summary and re-enqueues
affected callers. This feedback updates type-state relations and their
dependent results on the fixed call graph.



\subsection{Modeling Type-State Relations}
\label{sec:state-model}
\paragraph{Versioned state.}
An analysis state is $A=(\Gamma,H,\phi)$. The environment $\Gamma$ maps
scoped variable bindings to versioned value symbols. The heap $H$ records
the presence, value, type, and reference information of object locations.
The formula $\phi$ retains conditions and relations among symbols.
An instance attribute is identified by $\operatorname{Field}(o,a)$, where
$o$ is the object identity and $a$ the attribute name. Class attributes
have locations associated with their class objects, separate from instance
locations. Global bindings are qualified by their defining module.

Assignment creates a fresh value version at the updated binding or
location. A fact about the previous version remains a fact about that
value, while subsequent reads use the new version. Reference information
connects aliases to the same object identity. Presence distinguishes an
absent attribute, a present attribute whose value is \texttt{None}, and
unknown presence. Type terms include classes, subclass alternatives,
container types, literal types, symbolic types, unions, and an unknown
type. Semantic predicates describe relationships such as operand
compatibility and the type of an operation's result.

Alternative states retain complete relations. If an encoder choice is
associated with a particular field update, each alternative keeps that
association. This representation directly supports the distinction between
the illustrated call and its text-encoder alternative in
Section~\ref{sec:motivating-example}: the same field
location receives different value versions under different encoder cases.

\paragraph{Entries, events, and exits.}
A PSTC represents function $f$ by a set of relation cases, called path
contracts in Figure~\ref{fig:method-overview}:
\begin{equation}
\summary{f}=\{(p_i,m_i,q_i)\mid i\in I_f\}.
\label{eq:summary}
\end{equation}
Here $p_i$ describes the incoming state and the case's applicability
conditions. The ordered event graph $m_i$ retains guards, reads, writes,
operation requirements, calls, and exceptional behavior. The assertion
$q_i$ projects the normal exit onto the return value and caller-visible
state. Cases also carry source locations, read/write footprints,
dependency versions, coverage information, and evidence. A case ending
in failure or at an exploration boundary retains its event prefix and a
marker for the absence of a normal exit.

The figure's $P$, Path, and $Q$ fields correspond to $p_i$, $m_i$, and
$q_i$. In its first case, the incoming object \texttt{a} has type $A$ and
its attribute has type $B$; the exit relates the same object to an attribute
of type $C$ and a return value of type $T$. The Path field records the
conditions and operations establishing that change. As in the Hoare
triple introduced in Section~\ref{sec:hoare-background}, assertions connect
states across execution. Here the retained event graph also exposes the
checks and intermediate updates used for detection and refinement.

The distinction between applicability and operation safety is essential
for using the model. Entry matching selects a case applicable to the
incoming state. A later event $\operatorname{check}(R)$ tests an operation
requirement against the state immediately before that event and retains
both its successful and failing successors. Result facts become available
on the successful successor. Internal safety requirements therefore stay
at their execution positions rather than becoming assumptions used to
exclude an incoming call.

For example, in \texttt{y = 1 + "x"; if flag: ...}, addition fails before
the conditional is evaluated. Its failure prefix applies to both values
of \texttt{flag}. Preserving that prefix makes the failure available to
callers even though the function has no normal exit along this execution.
Similarly, an exceptional exit carries writes performed before the failed
operation.

For the encoder example, let $\ell=\operatorname{Field}(o,\mathtt{body})$
and let $e$ denote an exact instance of one of the shown encoder classes.
For an incoming string, the normal effects include
\begin{equation}
\begin{aligned}
\typeat{\ell_{\mathrm{in}}}&=\mathtt{str},\\
\typeat{e}=\mathtt{TextEncoder}&\Rightarrow
  \typeat{\ell_{\mathrm{out}}}=\mathtt{str},\\
\typeat{e}=\mathtt{Utf8Encoder}&\Rightarrow
  \typeat{\ell_{\mathrm{out}}}=\mathtt{bytes}.
\end{aligned}
\label{eq:encoder-effects}
\end{equation}
Both cases return \texttt{None}. Their event graphs retain the body guard,
method call, and write establishing the effect. A further case of \texttt{encode\_message} preserves a bytes body
when the helper's \texttt{isinstance(msg.body, str)} guard is false. The entry and exit assertions are
thus endpoints of an ordered behavior relation.

\subsection{Phase I: Context-free Code Simplification}
\label{sec:extraction}
\paragraph{Context-free Code Extraction.}
The figure's term \emph{context-free} means independent of the unknown
entry types and state. A computation belongs to this part, $C_I$, when
local semantics and available callee PSTCs determine its type behavior
without instantiating that unknown information. The remaining computations
form the context-dependent part $C_D$, labeled \emph{Context-sensitive Code
Snippet} in the figure. Dependence follows values, object state, guards,
and callee effects. Both parts retain source positions and control
dependencies for assembly in execution order.

For the schematic function \texttt{f(a)}, the literal assignment
\texttt{c = 1} establishes an integer input to \texttt{g}. The depicted
callee contract has an integer entry and integer return, allowing the
analyzer to determine the return type of \texttt{g(c)}. The subsequent
call \texttt{h(v, a)} still depends on the unknown input \texttt{a}.
The partition therefore separates the computation establishing
\texttt{v: int} from the use that requires additional type alternatives.
Its classification uses the callee's state dependencies and effects as
well as its parameter and return types.

\paragraph{Context-free Type Error Detection.}
The analyzer checks $C_I$ using local operation semantics and applicable
callee contracts. A successful operation supplies facts to its uses in
$C_D$; a conflict enters Phase~III. Each supplied fact records its value
or location version, condition, type relation, and supporting evidence.
In the figure, \texttt{v: int} is supplied at the use of \texttt{v} in
\texttt{h(v, a)}, allowing the computation of \texttt{g(c)} to be shared
across the candidate behaviors of \texttt{a}.

Entry-independent computations can contain branches, visible writes, and
exceptions. Their conditions, checks, effects, and exceptional successors
remain in the assembled relation. A result fact is supplied after its
producing operation succeeds, and facts affected by an intervening write
are updated at their later uses. Thus, simplification shares inference
while retaining the execution behavior that determines each fact's validity.

\subsection{Phase II: PSTC Generation}
\label{sec:pstc-generation}
\paragraph{Type Path Extraction.}
The analyzer explores explicit source paths in $C_D$ within the configured
bounds. Each template retains guards and operations in execution order,
together with applicable facts from Phase~I. The \emph{Type Paths} shown
in Figure~\ref{fig:method-overview} are these templates. A source path can
support several implicit type behaviors: \texttt{x + y}, for example,
can describe integer addition, string concatenation, or an operation
implemented by repository classes. Resolving these alternatives associates
each operand combination with its result and state effects.

\paragraph{Type Constraint Resolution.}
The LLM receives the explicit path, relevant use-site facts, repository
type information, and available callee PSTCs. It proposes implicit
input/output combinations and corresponding entry and exit assertions.
The analyzer assembles them with the ordered checks and effects from both
$C_I$ and $C_D$. Source paths determine the locations and ordering of the
operations to which each candidate relation applies. Candidate records
retain their inference context and source dependencies for subsequent
inspection. The resulting cases constitute the function's PSTC.

Inferred relations participate in propagation before conflict validation.
At an input state not covered by the available cases, the analyzer requests
supplementary analysis or retains an unresolved outcome. An operation or
annotation conflict enters Phase~III, which expands and re-infers the
relevant behavior in context. Uncovered inputs and conflicts therefore
focus additional construction on the relations needed by current callers.

\subsubsection{Call-site Composition and Operation Checks}
\label{sec:composition}
PSTC generation and application use the same call-site composition
procedure. At each call, the analyzer evaluates the call target and arguments in
source order and checks argument binding. It constructs a substitution
$\theta$ from formal parameters, object references, and environment
locations to the calling state. Internal symbols and newly allocated
objects are fresh for each instantiation. When two arguments refer to the
same actual object, the corresponding formal references retain that alias.

For a case $(p_i,m_i,q_i)$, applicability is evaluated under
\begin{equation}
\operatorname{Facts}(A)\land\theta(p_i).
\label{eq:applicability}
\end{equation}
The analyzer then processes $\theta(m_i)$ in event order. A guard refines
the current condition. A read obtains the current version of its location.
A write updates that location. An operation requirement is checked against
the state established by the preceding events. The exit assertion projects
the resulting return and shared state; it does not perform the updates a
second time.

A definite write to one location representing a single object creates a strong update.
For several possible targets, a weak update retains the possible old and
new contents. A previous location fact survives a call when that location
is provably disjoint from the possible write footprint. An unknown write
invalidates precise facts for the affected region. For example, after \texttt{a.prefix = b"x"} and
\texttt{b.prefix = "y"}, reading \texttt{a.prefix} yields a string
when \texttt{a is b}, but bytes when they are distinct ordinary objects.
The actual alias binding thus determines the state of a subsequent use.

For a requirement $R$ and the current facts $F$, the check examines the
success condition $F\land R$ and the failure condition $F\land\neg R$.
A supported failure alternative becomes a candidate conflict, associated
with its source location and state dependencies. Indeterminate conditions
remain unresolved. Subsequent events on a normal continuation use the
success facts; exceptional continuations carry the state accumulated up
to the failed event. Loop and recursive expansion follows the configured
budgets and records the remaining frontier.

For Figure~\ref{fig:motivating-example}, $\theta$ binds the helper's message
object to the caller's message and its encoder to \texttt{Utf8Encoder}.
The applicable case writes a bytes value to the body's new version.
The final addition is checked using that version, producing the bytes/string
conflict. The text-encoder case instead establishes a string version and
satisfies the same operation requirement.

\subsection{Phase III: Type Conflict Analysis}
\label{sec:conflict-analysis}
\paragraph{Callee PSTC Expanding.}
A candidate conflict links an operation or annotation check to the relations
supporting its judgment. Validation begins with those dependencies and
reconstructs the relevant behavior in the current calling context.
The expansion unit is a \emph{type path}: the ordered source behavior
associated with a relation case's entry and exit assertions. Those
assertions guide selection of the necessary callee paths on the fixed
call graph.

The analyzer follows the candidate's control, value, and alias dependencies
to select callee cases. It connects their type paths to the relevant caller
fragments, preserving conditions, operation requirements, writes, and
exceptional control flow. The resulting \emph{error trace} supplies both
the source behavior behind the disputed relation and the call-site-specific symbolic
bindings of the calling context. This error trace corresponds to the
figure's \emph{Conflict Trace}.

\paragraph{Conflict Path Analysis.}
The analyzer re-infers behavior over the expanded trace to determine
whether the conflict persists and whether the propagated relation
requires revision.

The figure's \emph{Missing Path Contract} is a relation supplied by this
analysis to supplement or correct the callee PSTC. Refinement can
supplement an entry case, separate alias alternatives,
restore a condition--effect association, or revise an inferred type
combination. A correction derived for a particular calling condition keeps
that condition. Revised summaries receive new dependency versions; results
that use an older version are invalidated and the affected caller entries
are re-enqueued. Analysis stops refining a candidate when its disposition
is established, no new evidence is obtained, or the budget is exhausted.
Algorithm~\ref{alg:analysis} summarizes this propagation and feedback.

\begin{algorithm}[t]
\caption{Propagation with conflict-driven relation refinement}
\label{alg:analysis}
\begin{algorithmic}[1]
\Require Repository $P$, fixed call graph $G$, entries $E$, budgets $B$
\Ensure Runtime reports, annotation-contract reports, and unresolved cases
\State Initialize summaries and a callee-first worklist $W$ from $G,E$
\While{$W\neq\emptyset$ and analysis budget remains}
  \State Remove a function entry $(f,A)$ from $W$
  \State Construct or supplement $\summary{f}$ using shared facts and type paths
  \State Compose relations under $A$ and collect operation and annotation conflicts
  \For{each candidate conflict $c$}
    \State Trace supporting conditions, values, aliases, and callee cases
    \State Expand relevant type paths in $G$ and re-infer in the calling context
    \If{a supporting relation changes}
      \State Revise its summary and invalidate dependent results
      \State Enqueue affected caller entries in $W$
    \EndIf
    \State Classify $c$ using its error category and current evidence
  \EndFor
\EndWhile
\State Retain remaining uncovered states and undecided candidates as unresolved
\end{algorithmic}
\end{algorithm}

\paragraph{Runtime type-error reports.}
A runtime report requires evidence that the expanded trace reaches an
operation with incompatible state and that its target protocol exception
escapes the declared entry. For a repository entry, the trace includes
the initialization and calls establishing the state; an open-entry report
records its symbolic arguments and state assumptions. A handler that
consumes the failure produces a handled-failure outcome. A corrected
relation that removes the conflict produces a dismissed candidate.
Insufficient evidence for either disposition produces an unresolved case.
Reports retain the expanded trace, supplementary context, re-inference
result, and exception disposition.

\paragraph{Annotation-contract reports.}
Existing annotations are checked against the requirements and effects
inferred from function behavior. For an input annotation, the analysis
examines whether the declared inputs are semantically compatible with the
inferred requirements. For return and state declarations, it examines
whether the inferred effects are compatible with the declared types.
A conflict enters the same dependency tracing and contextual re-inference
process. A confirmed report retains the declaration location, incompatible
relation, and supporting behavioral evidence. Its confirmation concerns
the declaration--behavior conflict and does not require a repository call
that raises a runtime exception.

\subsection{Reuse and Analysis Cost}
\label{sec:reuse}
Summary templates reuse extracted relations while freshening their internal
symbols at each application. Reuse of an already instantiated result
additionally matches the relevant entry state, aliases, call targets,
semantic configuration, assumptions, remaining expansion bounds, and
dependency versions. Cache keys include the transitive dependencies of
reads, checks, and effects. Instances with unknown read/write footprints
are recomputed. A cached failure is reconsidered under the current entry's
reachability and exception context.

The computational work comprises shared fact inference, path and type
alternative construction, composition, conflict expansion, contextual
re-inference, and dependency maintenance. Reuse acts at several of these
points: independent facts serve multiple alternatives, event prefixes and
callee templates serve multiple cases, and valid instances serve matching
calls. Refinement revisits affected callers through explicit dependencies.
Section~\ref{sec:rq2} measures the complete cost and compares these forms
of reuse with direct function-body analysis.
